% Original nonrelativistic limit worked out by component.
\subsection*{Worked calculation: normalize a Dirac spinor and take its low-energy limit}
The four-component spinor just introduced can be understood
by calculating its norm directly. Let
$\xi_s^\dagger\xi_{s'}=\delta_{ss'}$ and define
\begin{equation}
\chi_s=\frac{c\,\vect\sigma\cdot\vect p}{E+mc^2}\xi_s,
\qquad E=\sqrt{m^2c^4+c^2\vect p^{\,2}}>0.
\end{equation}
Because $(\vect\sigma\cdot\vect p)^\dagger
=\vect\sigma\cdot\vect p$ and
$(\vect\sigma\cdot\vect p)^2=\vect p^{\,2}I$,
\begin{align}
\chi_s^\dagger\chi_{s'}
&=\frac{c^2}{(E+mc^2)^2}
 \xi_s^\dagger(\vect\sigma\cdot\vect p)^2\xi_{s'}\notag\\
&=\frac{c^2\vect p^{\,2}}{(E+mc^2)^2}\delta_{ss'}.
\end{align}
The positive-energy dispersion relation implies
$c^2\vect p^{\,2}=E^2-m^2c^4
=(E-mc^2)(E+mc^2)$.
Thus
\begin{align}
\delta_{ss'}+\chi_s^\dagger\chi_{s'}
&=\left[1+\frac{E-mc^2}{E+mc^2}\right]\delta_{ss'}\notag\\
&=\frac{2E}{E+mc^2}\delta_{ss'}.
\end{align}
Multiply the stacked spinor $(\xi_s,\chi_s)^{\mathsf T}$
by $\sqrt{(E+mc^2)/(2E)}$. The normalization factor
cancels the previous fraction exactly:
$u_s^\dagger(\vect p)u_{s'}(\vect p)=\delta_{ss'}$.

\paragraph{Explicit expansion for small momentum.}
Introduce $\epsilon=|\vect p|/(mc)$ and assume
$\epsilon\ll1$. Expand the positive energy:
\begin{align}
E&=mc^2\sqrt{1+\epsilon^2}\notag\\
 &=mc^2\left(1+\frac{\epsilon^2}{2}
   -\frac{\epsilon^4}{8}+O(\epsilon^6)\right)\notag\\
 &=mc^2+\frac{\vect p^{\,2}}{2m}
   -\frac{\vect p^{\,4}}{8m^3c^2}
   +O\left(\frac{|\vect p|^6}{m^5c^4}\right).
\end{align}
Consequently
\begin{equation}
\frac{c\,\vect\sigma\cdot\vect p}{E+mc^2}
=\frac{\vect\sigma\cdot\vect p}{2mc}
+O(\epsilon^3),
\end{equation}
where the big-$O$ term is understood relative to the
dimensionless spinor components. The lower two components
are suppressed by one power of $|\vect p|/(mc)$.
The difference $E-mc^2$ becomes the nonrelativistic kinetic
energy at leading order. These two limits explain why
the two-component Pauli spinor is a reliable low-energy
description, together with the magnetic coupling derived
earlier in this chapter.

\paragraph{Checkpoint.}
For momentum along $+z$,
$\vect p=(0,0,p)$ and $\xi_\uparrow=(1,0)^{\mathsf T}$,
write the complete four-component $u_\uparrow$
and check all four squared moduli sum to one.
Then repeat for $\xi_\downarrow$ and verify orthogonality.
